
\[
j_*(y')+f^*(x)=0.
\]
Il
faut trouver \(y\in K^{\bullet}(Y)\) tel que:
\begin{equation}\tag{3.7.1}
x=-i_*(y);\qquad y'=g^*(y).\lambda_{-1}(F).
\end{equation}
Comme \(u\) admet un inverse à gauche \(u'\), on a nécessairement
\[
y=u'(y',x)=g_*(y').
\]
Il faut donc montrer que \(y\) vérifie les relations (3.7.1).

Pour montrer que \(x=-i_*(y)\), on part de la relation \(f^*(x)=-j_*(y')\) et on lui applique \(f_*\). On obtient donc: \(f_*(f^*(x))=-f_*(j_*(y'))=-i_*(g_*(y'))=-i_*(y)\) par définition de \(y\). Or d'après la proposition 3.6, \(f_*(f^*(x))=x\); d'où la relation voulue.

Pour achever la démonstration, il faut montrer que
\[
y'=g^*(g_*(y'))\lambda_{-1}(F).
\]
Posons \(z=y'-g^*(g_*(y'))\lambda_{-1}(F)\). On obtient: \(g_*(z)=g_*(y')-g_*(y')\,g_*(\lambda_{-1}(F))=0\). D'autre part, on a par hypothèse: \(j_*(y')=-f^*(x)\), d'où d'après ce qui précède \(j_*(y')=f^*(i_*(y))\), et d'après la proposition 3.4: \(j_*(y')=j_*\bigl(g^*(y)\lambda_{-1}(F)\bigr)\); par suite, on a également \(j_*(z)=0\). Or \(\Ker(j_*)\cap\Ker(g_*)=0\). En effet, si \(j_*(z)=0\), \(j^*(j_*(z))=0\); or \(j^*(j_*(z))=z(1-L)\) (prop. 2.7), et pour que \(z(1-L)=0\), il faut qu'il existe \(z'\in K^{\bullet}(Y)\) tel que \(z=z'\lambda_{-1}(F)\) (VI,5.10); en appliquant \(g_*\) ceci entraîne \(z'=g_*(z)=0\), donc \(z=0\).

Supposons maintenant \(X\) noethérien et \(Y\) régulier. Alors \(Y'\) est aussi régulier puisque c'est un fibré projectif sur \(Y\), et pour tout point \(y'\) de \(Y'\), l'anneau local \(\underline{O}_{X',y'}\) de \(X'\) en \(y'\) est un anneau régulier, car \(\underline{O}_{Y',y'}\) est quotient de \(\underline{O}_{X',y'}\) par un élément régulier (voir EGA \(0_{\mathrm{IV}}\) 17.1.8).

